\documentclass[11pt,a4paper]{article}
\usepackage[margin=25mm]{geometry}
\usepackage{amsmath,amssymb,amsthm,booktabs}
\newtheorem{lemma}{Lemma}\newtheorem{theorem}{Theorem}
\title{Five order 10 counts from two overgroups and $J_1$}
\author{Human Compression project}\date{}
\begin{document}\maketitle
Let $G=O'N$ and put
\[
 x_J(p,q)=|\{(r,s)\in J^2:|r|=p,\ |s|=q,\ |rs|=2\}|.
\]
Write $e_J(p,q)$ for the number of these pairs that generate $J$.
We recover five additional entries of the fixed Connor--Leemans
Table~1. The complete subgroup roster and acquired character tables
are borrowed premises, with their previous provenance retained.
Row numbers below refer to that roster. We reuse the accepted
order 20 argument, including the exact structure of its intersection
subgroup; no large group is enumerated in this application.

\begin{lemma}\label{unique}
Every element $s$ of order 10 belongs to a unique maximal subgroup
$A$ in row 7 and a unique maximal subgroup $B$ in row 11.
\end{lemma}
\begin{proof}
There is one cyclic subgroup class of order 10, row 549, and it
occurs in cyclic order 20, row 493. Thus every such $s$ is contained
in a cyclic subgroup of order 20. The primary ATLAS table gives
$|C_G(s)|=20$ for the sole element class 10A. Consequently $C_G(s)$
is that cyclic group, and has a unique involution. The row 7 maximal
subgroups are involution centralizers, so exactly one contains $s$.
The unique row 11 maximal subgroup follows from the already proved
uniqueness of the row 11 overgroup of $\langle s^2\rangle$, of order
five. That uniqueness was obtained by the direct incidence
$17778376\cdot36=640021536$.
\end{proof}

Let $H$ denote a row 32 subgroup, of order 2880. The accepted order
20 proof gives a normal subgroup $N\cong A_6$ with $H/N\cong Q_8$.
The preimage $K$ of the center of $H/N$ is $A_6\times C_2$.
Each $H$ belongs to exactly one row 7 and one row 11 maximal subgroup,
and is their intersection. There are 160005384 conjugates of $H$.

\begin{lemma}\label{overlap}
For $(p,q)=(3,10),(4,10),(6,10),(10,12)$, a pair of the required
orders contained in both a row 7 and a row 11 maximal subgroup
belongs to their row 32 intersection. Its intrinsic counts are
\[
 x_H(3,10)=1440,\quad x_H(4,10)=2880,\quad
 x_H(6,10)=1440,\quad x_H(10,12)=0.
\]
\end{lemma}
\begin{proof}
Intersecting the complete descendant lists of rows 7 and 11 with
the complete ancestor lists of the two cyclic orders leaves exactly
\[
\begin{array}{c|l}
(p,q)&\text{possible generated subgroup rows}\\\hline
(3,10),(6,10)&32,57,60,61,104,210,282\\
(4,10)&32,57,60,61,104,210,312,401,402,403,490,493\\
(10,12)&32,57,210.
\end{array}
\]
Every row in these lists occurs below row 32. Conjugacy therefore
supplies an actual row 32 subgroup containing the generated subgroup
of the pair. Its unique maximal overgroups must be the two already
determined by Lemma~\ref{unique}; it is their intersection. This
uses existence from the complete subgroup diagram, not products of
cover multiplicities.

If $|s|=10$ in $H$, its image in $Q_8$ is the unique involution:
it has order dividing two and cannot be the identity, since $A_6$
has no element of order ten. If $rs$ is an involution, its image is
also central. Hence $r,s\in K=A_6\times C_2$. In this group write
$s=(b,-1)$ with $|b|=5$. For $|r|=3$, one has $r=(a,1)$ with
$|a|=3$; for $|r|=6$, one has $r=(a,-1)$ with $|a|=3$.
For $|r|=4$, the first coordinate has order four and either sign
is possible. The accepted literal values
$x_{A_6}(3,5)=x_{A_6}(4,5)=1440$ give the first three counts.
There is no element of order 12 in $A_6\times C_2$, giving the last.
A separate literal enumeration of the complete 720-element direct
product verifies all four values.
\end{proof}

\begin{lemma}\label{j1}
Put $J=J_1$. Then
\[
\begin{aligned}
e_J(3,10)&=x_J(3,10)-1463\cdot120,\\
e_J(6,10)&=x_J(6,10)-1463\cdot120-2926\cdot120,\\
e_J(10,15)&=x_J(10,15)-2926\cdot24.
\end{aligned}
\]
Every nongenerating pair in these three cases is contained in a
row 7 or a row 11 maximal subgroup of $G$.
\end{lemma}
\begin{proof}
The complete list of maximal subgroups of $J$, with the number of
actual subgroups in each class, is
\[
\begin{array}{c|rrrrrrr}
\text{row}&119&237&277&285&286&366&397\\
\text{order}&660&168&120&114&110&60&42\\
\text{number}&266&1045&1463&1540&1596&2926&4180.
\end{array}
\]
Row 119 is $L_2(11)$, whose nonidentity element orders divide 5 or 6,
or equal 11, and so it contains no element of order 10. The other
orders exclude our pairs except for rows 277 and 366, which the
roster identifies unambiguously as $A_5\times C_2$ and
$D_{10}\times D_6$. Here $D_{2n}$ has order $2n$.

The exact small product calculations are
\[
\begin{array}{c|rrr}
 &(3,10)&(6,10)&(10,15)\\\hline
A_5\times C_2&120\ [120]&120\ [120]&0\\
D_{10}\times D_6&0&120\ [60]&24\ [30].
\end{array}
\]
The bracketed number is the order of the subgroup generated by
every counted pair. These statements have literal certificates:
$A_5$ is the full set of 60 even permutations of five letters;
$D_{2n}$ is the set of affine permutations $i\mapsto a\pm i$ of
$\mathbb Z/n\mathbb Z$. Disjoint permutation actions realize the
two products. The reader checks every pair, its three orders, and
the complete closure under its two generators.

For the last entry all 24 pairs generate the same normal subgroup
$C_5\times D_6$ of the displayed $D_{10}\times D_6$. This is row 465.
Its global class length is 7680258432, so its normalizer in $G$ has
order 60. Each row 366 subgroup normalizes its unique row 465
subgroup and has that order. Two row 366 subgroups therefore cannot
have the same row 465 subgroup. In the other nonzero entries the
pairs generate the entire indicated maximal subgroup. Thus no
pair is counted twice among these proper contributions, either
within a type or between the distinct generated orders. Subtracting
them from $x_J$ proves the three identities.

Finally, $A_5\times C_2$ has a central involution and is contained
in its row 7 centralizer. The group $D_{10}\times D_6$ normalizes
its cyclic rotation subgroup $C_5$. Uniqueness of the row 11
maximal overgroup of that $C_5$ implies that its normalizer
normalizes that maximal subgroup. A proper maximal subgroup of
the simple group $G$ is self-normalizing. Hence
$D_{10}\times D_6$ is contained in the row 11 maximal subgroup.
This covers every proper contribution just counted.
\end{proof}

\begin{theorem}
The five ordered automorphism-orbit counts are
\[
\begin{array}{c|rrrrr}
(p,q)&(3,10)&(4,10)&(6,10)&(10,12)&(10,15)\\\hline
 e_G(p,q)/|\operatorname{Aut}(G)|&7&285&953&1969&3072.
\end{array}
\]
\end{theorem}
\begin{proof}
The acquired complete maximal character tables have nonzero pair
counts only in rows 6, 7 and 11 for $(3,10),(6,10)$; only in rows 7
and 11 for $(4,10),(10,12)$; and only in rows 6 and 11 for $(10,15)$.
Row 6 is $J_1$. All its proper contributions are already covered by
rows 7 or 11, by Lemma~\ref{j1}. Its generating pairs, however,
cannot lie in any other maximal subgroup of $G$.

We may therefore replace the contribution $[G:J]x_J$ in the maximal
sum by $[G:J]e_J$. For rows 7 and 11, each pair is counted at most
once in each class, by Lemma~\ref{unique}. Their double counting
is exactly $160005384x_H$, by Lemma~\ref{overlap}. In the last case
the row 7 contribution is zero, so there is no overlap term.
The resulting generating-pair count is thus
\[
 e_G=x_G-\sum_M[G:M]x_M
     +[G:J](x_J-e_J)+160005384x_H,
\]
with the $J$ term set to zero in the two cases where $x_J=0$, and
the $H$ term zero for $(10,15)$. The acquired character evaluations,
the three identities of Lemma~\ref{j1}, and the four literal overlap
values give the displayed answers. Division is by
$|\operatorname{Aut}(G)|=2|G|$, whose action on generating pairs is
free. No division by duality is made. All five answers agree with
the fixed printed entries.
\end{proof}

\paragraph{Attribution and limits.}
The numerical assertions are from Connor and Leemans,
\emph{Algorithmic enumeration of regular maps}, Table~1. The complete
roster, named small subgroup types, ATLAS centralizer data and
character evaluations remain explicit borrowed premises. The new
argument separates generating pairs of $J_1$ from its proper
contributions and reuses the accepted order 20 intersection theorem.
Its finite certificates concern only the displayed small products.
This component does not settle the remaining chirality assertions,
claim priority for its elementary group theory, or complete the
original manuscript.
\end{document}
